\begin{thebibliography}{99}
\footnotesize
\setlength{\itemsep}{1.5pt}

\bibitem{bmbl6}
US Centers for Disease Control and Prevention and National Institutes of Health.
\textit{Biosafety in Microbiological and Biomedical Laboratories}, 6th ed.
\url{https://www.cdc.gov/labs/bmbl/index.html} (accessed 29 October 2026).

\bibitem{wholbm4}
World Health Organization. \textit{Laboratory Biosafety Manual}, 4th ed. Geneva: WHO;
2020. \url{https://www.who.int/publications/i/item/9789240011311} (accessed 29 October
2026).

\bibitem{iso35001}
International Organization for Standardization. \textit{ISO 35001:2019, Biorisk management
for laboratories and other related organisations}. Geneva: ISO; 2019.

\bibitem{cms2026}
Centers for Medicare \& Medicaid Services, Division of Clinical Laboratory Improvement and
Quality. \textit{CLIA Laboratory Statistics, September 2026}.
\url{https://www.cms.gov/files/document/laboratory-statistics.pdf} (accessed 29 October
2026).

\bibitem{ravnholdt2025}
Ravnholdt AR, Ratnesar-Shumate SA, Santarpia JL. A review of accidental aerosol generation
in laboratories and laboratory-associated infections. \textit{Appl Biosaf}.
2025;30(3):229--238. doi:10.1089/apb.2024.0014.

\bibitem{dubuis2021}
Dubuis M-E, Duchaine C. Aerosol production during blood and urine pre-analytical
processing and handling in a hospital biochemistry clinical laboratory during the COVID-19
pandemic. \textit{Front Public Health}. 2021;9:643724. doi:10.3389/fpubh.2021.643724.

\bibitem{liu2023}
Liu Z, Lv J, Zhang Z, Ma J, Song Y, Wu M, Cao G, Guo J. Three experimental common
high-risk procedures: emission characteristics identification and source intensity
estimation in biosafety laboratory. \textit{Int J Environ Res Public Health}.
2023;20(5):4479. doi:10.3390/ijerph20054479.

\bibitem{banada2010}
Banada PP, Sivasubramani SK, Blakemore R, Boehme C, Perkins MD, Fennelly K, Alland D.
Containment of bioaerosol infection risk by the Xpert MTB/RIF assay and its applicability
to point-of-care settings. \textit{J Clin Microbiol}. 2010;48(10):3551--3557.
doi:10.1128/jcm.01053-10.

\bibitem{ackelsberg2020}
Ackelsberg J, Liddicoat A, Burke T, Szymczak WA, Levi MH, Ostrowsky B, et al.
\textit{Brucella} exposure risk events in 10 clinical laboratories, New York City, USA,
2015 to 2017. \textit{J Clin Microbiol}. 2020;58(2):e01096-19. doi:10.1128/jcm.01096-19.

\bibitem{mmwr2002}
Centers for Disease Control and Prevention. Laboratory-acquired meningococcal disease ---
United States, 2000. \textit{MMWR Morb Mortal Wkly Rep}. 2002;51(7):141--144.

\bibitem{sejvar2005}
Sejvar JJ, Johnson D, Popovic T, Miller JM, Downes F, Somsel P, Weyant R, Stephens DS,
Perkins BA, Rosenstein NE. Assessing the risk of laboratory-acquired meningococcal
disease. \textit{J Clin Microbiol}. 2005;43(9):4811--4814.
doi:10.1128/jcm.43.9.4811-4814.2005.

\bibitem{mmwr2014}
Sheets CD, Harriman K, Zipprich J, Louie JK, Probert WS, Horowitz M, Prudhomme JC, Gold D,
Mayer L. Fatal meningococcal disease in a laboratory worker --- California, 2012.
\textit{MMWR Morb Mortal Wkly Rep}. 2014;63(35):770--772.

\bibitem{asmbruc2025}
American Society for Microbiology. \textit{Sentinel Level Clinical Laboratory Guidelines
for Suspected Agents of Bioterrorism and Emerging Infectious Diseases: Brucella species}.
2025. \url{https://asm.org/getmedia/ffb17144-8e7c-44b2-942d-9c4b8b557abf/brucella-guidelines_august2025.pdf}
(accessed 29 October 2026).

\bibitem{asmsentinel}
American Society for Microbiology. \textit{Laboratory Response Network (LRN) Sentinel Level
Clinical Laboratory Protocols}.
\url{https://asm.org/articles/cphmc/laboratory-response-network-lrn-sentinel-level-c}
(accessed 29 October 2026).

\bibitem{wholbm4bsc}
World Health Organization. \textit{Laboratory Biosafety Manual, 4th ed.: Biological Safety
Cabinets and Other Primary Containment Devices}. Geneva: WHO; 2020.
\url{https://www.who.int/publications/i/item/9789240011335} (accessed 29 October 2026).

\end{thebibliography}
